\begin{center}
  \textbf{\large 摘\quad 要}
\end{center}

\vspace{-0.5em}

本文复现了文档级事件因果识别（Event Causality Identification, ECI）方法
\textbf{CogECI}，并在 EventStoryLine v0.9（ESC）数据集上完成 5 折交叉验证。
CogECI 在 iLIF 的迭代式因果图骨架之上，引入\emph{句子图模块}（Sentence Graph Module, SGM）、
\emph{双流因果推理模块}（Dual-Stream Causal Reasoning Module, DSCRM）以及
\emph{上下文对齐损失}（Context Grounding Loss, CGL），
其第一阶段依赖大语言模型（LLM）为每个因果正例事件对标注"支撑该因果关系的句子集合"。

本项目的主要工作与结论如下。

\begin{enumerate}
  \item \textbf{完整打通了官方仓库的复现链路。}官方代码仓发布时缺失 \texttt{utils.py}、
        特征缓存 \texttt{*.pkl}、伪标签 CSV 与 BERT 权重共四类文件，其中 \texttt{utils.py}
        被 \texttt{main.py} 直接 \texttt{import}，缺失即无法启动。本文补齐上述依赖并逐一记录来源。
  \item \textbf{确认了评测口径。}通过标签分布实测（\texttt{exist} 版 pkl 正例全为 1）、
        正例总数交叉验证（4\,933 + 545 = 5\,478，与 iLIF 论文预处理后的正例数完全一致）
        与超参默认值三条独立证据，确认默认配置对应\textbf{因果存在性二分类}口径。
  \item \textbf{给出了 LLM 打标的可离线替代方案并做了定量检验。}受平台免费额度耗尽所限，
        Stage 1 未按论文原方案完成全量 LLM 打标。本文设计基于事件所在句的\emph{规则法}伪标签，
        并用已完成的 545 条 Qwen3.8-27B 标注做一致性对照：规则集相对 LLM 选句的
        \textbf{句级精确率 0.985}、F1 0.853，且 \textbf{97.4\%} 的样本中 LLM 所选句子恰好包含
        两个事件所在句，说明规则法捕捉到了 LLM 选句的核心，属"高精度、稍低召回"的保守替代。
  \item \textbf{完成了工程优化并严格验证了数值等价性。}原始实现在单文档上仅约 8\% GPU 利用率，
        5 折 $\times$ 30 轮外推需约 47 小时。本文对 5 处热点做了向量化改写，
        实测加速 \textbf{12--17 倍}，且损失值偏差为 0、预测概率偏差 $<1.2\times10^{-7}$、
        梯度范数偏差 $<10^{-8}$、图边选取结果逐元素相等。同时引入梯度检查点，
        将峰值显存由 6\,824\,MB 降至 2\,483\,MB，使 6\,GB 显存笔记本可完成全量训练。
  \item \textbf{主实验结果。}在因果存在性口径下，5 折平均
        $P=0.6812$、$R=0.6225$、$\bm{F_1=0.6493}$，
        其中句内 $F_1=0.7745$、跨句 $F_1=0.5882$。
        相较 iLIF 论文在 ESC 上报告的存在性 $F_1=62.5$（句内 71.2 / 跨句 59.1），
        整体提升 2.4 个点，且增益主要来自句内场景。
\end{enumerate}

\vspace{0.5em}
\noindent\textbf{关键词：}事件因果识别；文档级；图神经网络；上下文对齐；伪标签；复现研究

\vspace{1.2em}
\begin{center}
  \rule{0.85\textwidth}{0.4pt}
\end{center}
\vspace{0.3em}

\begin{center}
  \textbf{\large Abstract}
\end{center}

\vspace{-0.5em}
\noindent
This report reproduces \textbf{CogECI}, a document-level event causality identification
method, and evaluates it on EventStoryLine v0.9 with 5-fold cross-validation.
CogECI augments the iterative causal-graph backbone of iLIF with a Sentence Graph Module (SGM),
a Dual-Stream Causal Reasoning Module (DSCRM), and a Context Grounding Loss (CGL);
its first stage relies on an LLM to label, for each positive event pair, the set of sentences
that ground the causal relation.

We first repair the released repository (four categories of files are missing, including
\texttt{utils.py} which is imported by \texttt{main.py}), and confirm the evaluation protocol
to be \emph{causality existence} binary classification via three independent lines of evidence.
Due to exhaustion of the LLM free quota, we replace Stage-1 LLM labelling with a rule-based
offline scheme, and validate it against 545 Qwen3.8-27B annotations: sentence-level precision
0.985 and F1 0.853, with 97.4\% of LLM answers containing exactly the two event-bearing sentences.
We further vectorize five hotspots of the original implementation, achieving a 12--17$\times$
speed-up with numerically identical behaviour (loss difference 0, prediction difference
$<1.2\times10^{-7}$), and cut peak VRAM from 6\,824\,MB to 2\,483\,MB via gradient checkpointing.
The final 5-fold result is $P=0.6812$, $R=0.6225$, $F_1=0.6493$,
improving over the iLIF existence baseline of 62.5 on ESC.

\vspace{0.6em}
\noindent\textbf{Keywords:} event causality identification; document-level; graph neural network;
context grounding; pseudo labelling; reproduction study

\newpage
